% Glyph-free reconstruction after the first registered branch: rounds v2/v3, v4 design, criteria amendment, v4.1 and the sealed
% cross-generator test (numbers: numbers_g4.tex, generated by tools/make_numbers_g4.py). Part of Sec. D of the supplement.
% generated by tools/make_numbers_g4.py -- do not edit by hand
\begin{table*}[t]
\centering\small\setlength{\tabcolsep}{4pt}
\caption{All rounds of the glyph-free branch. Each verdict was taken by criteria fixed before that round's test run and none was revised afterwards. Readability: branch vs.\ comparator on the round's test data (first reader). D4 is development data of v4 and v4.1 and is not a test.}
\label{tab:g_rounds}
\resizebox{\linewidth}{!}{%
\begin{tabular}{lp{4.6cm}p{4.4cm}p{3.3cm}p{4.4cm}}
\toprule
Round & What changed & Test data & Readability & Registered verdict\\
\midrule
v1 & carrier-only layout, SIFT tracking & \lthreeALoc{} synthetic glyph-free clips (benchmark); \lthreeLtwo{} pseudo-text clips & \lthreeACarRead{} vs.\ \lthreeABfbRead{} (fallback); \lthreeBCarRead{} vs.\ \lthreeBBoneRead{} (\method) & not passed: \lthreeAPassed{} of \lthreeACriteria{} criteria met (geometry and stability missed)\\
v2 & motion-model selection; in-carrier layout with a font search & 449 synthetic clips (benchmark); 260 CogVideoX pseudo-text clips & 66.1 vs.\ 81.3; 55.4 vs.\ 78.1 & not passed: readability, IoU, height ratio 1.27, jitter; the font search chose display and calligraphic faces\\
v3 & v2 tracking; frozen default fonts & 57 locatable CogVideoX synthetic clips; the benchmark's 71 pseudo-text clips & 89.5 vs.\ 75.4; 98.6 vs.\ 90.1 & not passed: synthetic 6 of 8 criteria failed; pseudo-text only BG (40.67\,dB $<$ 41\,dB)\\
v4 (v4-LT) & dense-flow tracker with keyframe anchors; layout in carrier $\cap$ allowed region & D4 (development, two generators) & -- & not a test; on D4 two implementation violations of the edit-support constraint were found\\
v4.1 & unified region compositor; inward feather; outside-region invariant & sealed T4: 265 synthetic and 631 pseudo-text clips, two generators & pooled +20.0 pp [+14.8, +25.5] over the fallback & not passed: 3 of 34 gates failed (CogVideoX height ratio and normalised jitter; Wan2.2 normalised flicker)\\
\bottomrule
\end{tabular}}
\end{table*}


\subsection{Rounds v2 and v3}
\label{sec:supp_g_rounds}
After the first registered branch failed (above), two more rounds were registered on 2026-10-24, each before its own test run; \cref{tab:g_rounds} lists every round with its verdict. \textbf{v2} replaced the per-frame tracking by one choice per clip among a static, a homography and a dense-flow trajectory, each smoothed jointly over the clip and selected by its held-out photometric error with a 10\% margin over the static one, and laid the text out inside the carrier with a font search that maximised the text height. On the \gVtwoSynN{} locatable synthetic clips of the benchmark it read \gVtwoSynRead\% against \gVtwoSynFb\% for the static-hint fallback (v1 in the same run: \gVoneSynRead\%). The tracking change alone kept the readability (\gVtwoAoneRead\%) and the layout change alone lost it (\gVtwoAtwoRead\%): the font search chose condensed, display and calligraphic faces that the reader could not read and enlarged the text (median height ratio \gVtwoHeight). On \gVtwoNatN{} pseudo-text clips generated by CogVideoX-5B it read \gVtwoNatRead\% against \gVtwoNatBone\% for \method. \textbf{v3} kept the tracking and returned to the frozen default fonts. Because its design drew on the v2 diagnostics, it was tested on two sets it had never been run on: \gVthreeSynN{} locatable synthetic clips built from CogVideoX-5B clips (of \gVthreeSynAll) and the \lthreeLtwo{} pseudo-text clips of the benchmark. It read \gVthreeSynRead\% against \gVthreeSynFb\% (\gVthreeWins{} vs.\ \gVthreeLosses{} discordant clips, $p=\gVthreeP$) but failed \gVthreeFailed{} of \gVthreeCriteria{} synthetic criteria, with jitter \gVthreeJit{} [\gVthreeJitLo, \gVthreeJitHi] above the fallback; on the pseudo-text clips it read \gVthreeLtwoRead\% against \gVthreeLtwoBone\% ($p=\gVthreeLtwoP$) and failed only the background criterion (\gVthreeLtwoBg\,dB against a limit of 41\,dB), because a layout that follows the carrier can extend beyond the region outside which the background is measured. Both verdicts stand.

Diagnostics run while designing v4---on the v1--v3 outputs, the CogVideoX synthetic set and the development set D4 below, never on the sealed set---explained three recurring failures. (i)~The registered jitter and flicker scale with the edge energy inside the text box: the branch writes high-contrast text where the fallback often writes faint or no text, so their raw difference mostly measures rendering contrast; normalised by edge energy, no branch version was less stable than the fallback. (ii)~The hint box describes frame~0, but every version had used it as the text position on the middle frame, which displaces the text on moving clips by the motion between the two frames. (iii)~IoU with the old text's footprint favours the fallback whenever the hint box is derived from that footprint. These findings changed the metrics and the method of v4, never an earlier verdict.

\subsection{v4: constrained carrier-space reconstruction}
\label{sec:supp_g_v4}
\textbf{Goal and data.} v4 was registered on 2026-10-24 with the requirement that correctness, geometric stability and edit locality improve together, and with the rule that the branch enters \method only if it passes a sealed test. Prompts were generated once: the development set D4 uses 20 target strings and 3 scene templates, the sealed set T4 another 20 strings and 3 templates, each string in four of the seven materials (80 prompts per set); neither set shares a string or a template with the other or with the benchmark's 40 strings and 5 templates. Every prompt was rendered by CogVideoX-5B (50 steps, guidance 6, 720$\times$480, 49 frames at 8\,fps) and Wan2.2-TI2V-5B (40 steps, guidance 5, 832$\times$480, 49 frames at 24\,fps), with fixed seeds: 2 per prompt for D4 and 6 for T4. T4 was raised from 2 to 6 seeds before any of its clips was classified, because 2 seeds would have left fewer than the registered 60 synthetic clips per generator. The frozen localiser classifies every clip; clips with readable matches are erased as in v1 to give synthetic glyph-free clips with geometry truth, pseudo-text clips form a natural subset that is judged separately, and clips without a detectable text region are only counted.

\textbf{Method.} The inputs are the clip, the frame-0 hint box and the target, as in v1. \emph{Tracking:} each step is estimated by dense optical flow inside the carrier region with a per-point forward--backward check (2\,px), at least 30\% consistent points and 30 inliers; sequential RANSAC fits up to two motion layers and keeps the one with the lower photometric error on the carrier core, provided it beats the identity by a factor of 1.15. Keyframes are re-localised against the pre-warped reference to anchor drift, sparse SIFT matches replace the flow where it fails, and the trajectory is fused from anchors, bidirectional chains, interpolation and hold. Each clip chooses a smoothing window of 3, 7 or 13 frames by photometric error (7 unless another is at least 1\% better), and is tracked rather than kept static only with at least three anchors, coverage $\geq 0.5$, an excursion of at least 1\,px and a photometric margin of 0.10. The frame-0 hint is mapped to the reference (middle) frame through the estimated motion. \emph{Layout:} the feasible region is the carrier's extent (scanned as in v1, inset by 0.08 text heights) intersected with the allowed region---the hint box dilated by the larger of 4\,px and half its height, i.e.\ the region outside which background PSNR is measured---back-projected to the reference frame through the trajectory. One to three lines are chosen by matching the text box's aspect ratio to the hint's; a single line is at most the hint's height and 0.75 of the usable height, several lines at most 0.85 of it, and the width at most 0.92; text that would fall below the larger of 12\,px and 0.2 hint heights falls back to \method. Fonts are the frozen defaults of v3; ink and erasure are as in v1.

\textbf{Design.} Every design choice used only the v1--v3 outputs, the CogVideoX synthetic diagnostic set and D4. Dense flow replaced sparse feature chains because generated carriers are smooth (on a quarter of the moving synthetic clips SIFT found no reliable feature); the hint mapping removed the displacement of (ii); the per-clip smoothing reconciled synthetic clips, whose erased region moves with smoothed truth geometry, with natural ones; and the readability floor was lowered from 0.35 to 0.2 hint heights after it had rejected readable text on large hints. T4 was not opened during the design.

\subsection{Acceptance criteria and their amendment}
\label{sec:supp_g_criteria}
The criteria registered with v4 were amended once, after the D4 calibration and before any T4 clip was classified or inspected. The amendment is labelled a \emph{prospective amendment after D4 calibration}, not part of the original registration; it changes no earlier verdict, and D4 counts as development evidence only. It made four changes. (1)~\emph{Correctness:} superiority is tested on the two generators pooled with equal weights (CI lower bound above zero and at least $+5$ points), with non-inferiority per generator (CI lower bound $\geq -5$ points), because the static fallback alone read 93\% on the Wan2.2 part of D4 and left no room for a per-generator superiority test; the pooled readability must also reach v3's. On natural pseudo-text the branch must be non-inferior to the deployed \method per generator (CI lower bound $\geq -5$ points). (2)~\emph{Background:} relative to an identity re-encoding through the same decoder, encoder and region, per generator and subset (no-op $-0.3$\,dB), because even an identity re-encoding stays below 41\,dB on Wan2.2 clips; the paired differences to the fallback and to the deployed \method are reported. (3)~\emph{Tracking:} a controlled experiment with identical fonts, layout, erased background, colour and compositing that differs only in the trajectory (new tracker vs.\ the trajectory selection of v2/v3), judged on natural pseudo-text, because the erased regions of synthetic clips move with smoothed truth geometry and favour smooth trajectories. (4)~\emph{Containment:} mean containment of the text in the carrier of at least 0.90 and non-inferior to the fallback (CI lower bound $\geq -0.03$), with the original 95\% kept as a reported control, plus a separate severe-failure gate (human-confirmed, at most 5\% per generator) and a layout-acceptability gate. Temporal stability is judged by edge-normalised jitter and flicker relative to the fallback (CI upper bound $\leq 0$ per generator, i.e.\ per frame rate), for reason (i); raw values are reported. Geometry requires lower centre offset, drift and motion error than the fallback (CI upper bound below zero) and a median text-height ratio in $[0.8, 1.2]$; IoU and corner error are reported only. All CIs come from a bootstrap that resamples prompts within each generator (10{,}000 resamples, seed 0, percentile intervals), because the seeds of one prompt are correlated; synthetic and natural subsets are never pooled.

\subsection{v4.1: two locality fixes before the sealed test}
\label{sec:supp_g_v41}
On D4 the frozen v4 failed two automatic gates on CogVideoX: the background of the natural subset (\gDCogNatBgVfour{} against \gDCogNatBgNoop\,dB for the identity re-encoding) and the median height ratio (\gDCogHeight). The background failure came from two implementation violations of the registered rule that every edit, including antialiasing and erasure, stays inside the allowed region: the fallback copied the deployed system's whole output, which also changes pixels outside the region (\gCCogLeakFrac\% of the outside pixels and \gCCogLeakShare\% of the total change on the CogVideoX fallback clips), and the clipping feather reached about 3\,px beyond the region. v4.1 corrects exactly these two points. One compositor serves the branch and the fallback: outside the allowed region the output equals the decoded input, the feather runs inward only (opacity given by the exact Euclidean distance to the region boundary over 3\,px), and every frame is compared with the input outside the region before encoding, aborting the job on any difference. Fonts, sizes, layout, tracking, photometric models, smoothing, routing, fallback selection and every threshold are unchanged, and nothing was based on T4. The record states: \emph{``Version v4.1 corrects two implementation violations of the previously specified edit-support constraint, identified during D4 calibration. All changes were completed before inspecting T4 outcomes. The original v4-LT development results and all evaluation thresholds remain unchanged.''} On D4 the outside-region check passed on all \gCCogChecked{} CogVideoX and \gCWanChecked{} Wan2.2 outputs, the layouts were identical to those of v4, the background of the fallback clips relative to the identity re-encoding moved from \gCCogFbBgBefore{} to \gCCogFbBgAfter\,dB (CogVideoX, \gCCogFbN{} clips) and from \gCWanFbBgBefore{} to \gCWanFbBgAfter\,dB (Wan2.2, \gCWanFbN{} clips) with unchanged readability, and no output of the branch itself changed its reading. The natural-subset background then passed on CogVideoX (\gDCogNatBgVfourone\,dB) and the pooled development readability gain was \gDReadGain{} points [\gDReadGainLo, \gDReadGainHi]; the height ratio stayed at \gDCogHeight{} and was not tuned, so v4.1 entered the sealed test with one known automatic failure on D4. The code was frozen with a SHA-256 manifest; the v4 code and its D4 results are kept.

% generated by tools/make_numbers_g4.py -- do not edit by hand
\begin{table*}[t]
\centering\small\setlength{\tabcolsep}{4pt}
\caption{Sealed cross-generator test of the glyph-free branch v4.1: every registered gate, per generator (one run of the frozen code; prompt-clustered 95\% bootstrap CIs; differences are v4.1 minus the frozen code's static-hint fallback on the same erased inputs, or minus the deployed \method on natural pseudo-text). Synthetic: 109 (CogVideoX-5B, 36 prompts) and 156 (Wan2.2-TI2V-5B, 64 prompts) clips. The pooled primary gate (equal-weight mean of the two generators) is +20.0 points [+14.8, +25.5] (passed). A1: the trajectory selection of v2/v3. \ding{55}: failed gate. The two human gates could no longer change the verdict and were replaced by a diagnostic audit (Sec.~\ref{sec:supp_g4_audit}).}
\label{tab:g4}
\resizebox{\linewidth}{!}{%
\begin{tabular}{llrcrc}
\toprule
Gate & Threshold & CogVideoX-5B (8\,fps) & & Wan2.2-TI2V-5B (24\,fps) & \\
\midrule
\multicolumn{6}{l}{\textit{Correctness (synthetic glyph-free; natural pseudo-text is a separate task)}}\\
Readability gain over the static fallback, pp & CI lo $\geq -5$ & +26.6 [+18.1, +36.0] & \ding{51} & +13.5 [+7.5, +19.6] & \ding{51} \\
Natural pseudo-text: gain over deployed \method, pp & CI lo $\geq -5$ & +9.5 [+5.4, +13.8] & \ding{51} & +15.7 [+10.1, +21.5] & \ding{51} \\
\multicolumn{6}{l}{\textit{Temporal stability and locality}}\\
Edge-normalised jitter, difference to fallback & CI hi $\leq 0$ & \ensuremath{-}0.032 [\ensuremath{-}0.072, +0.014] & \textbf{\ding{55}} & \ensuremath{-}0.024 [\ensuremath{-}0.041, \ensuremath{-}0.008] & \ding{51} \\
Edge-normalised flicker, difference to fallback & CI hi $\leq 0$ & \ensuremath{-}0.069 [\ensuremath{-}0.113, \ensuremath{-}0.037] & \ding{51} & \ensuremath{-}0.001 [\ensuremath{-}0.033, +0.055] & \textbf{\ding{55}} \\
BG PSNR synthetic / identity re-encoding, dB & $\geq$ no-op $-0.3$ & 41.38 / 41.55 & \ding{51} & 40.77 / 40.87 & \ding{51} \\
BG PSNR natural / identity re-encoding, dB & $\geq$ no-op $-0.3$ & 41.26 / 41.37 & \ding{51} & 40.43 / 40.51 & \ding{51} \\
Spill outside the allowed region, \% & $\leq 2$ & 0.0 & \ding{51} & 0.0 & \ding{51} \\
Frame overflow, \% (fallback) & $\leq$ fallback & 0.0 (10.1) & \ding{51} & 1.3 (15.4) & \ding{51} \\
\multicolumn{6}{l}{\textit{Tracking and geometry}}\\
Same rendering, new tracker $-$ A1: raw jitter & $\leq +0.2$ & \ensuremath{-}0.52 [\ensuremath{-}0.92, \ensuremath{-}0.20] & \ding{51} & \ensuremath{-}2.04 [\ensuremath{-}2.56, \ensuremath{-}1.56] & \ding{51} \\
Same rendering, new tracker $-$ A1: raw flicker & $\leq +0.05$ & \ensuremath{-}0.080 [\ensuremath{-}0.143, \ensuremath{-}0.031] & \ding{51} & \ensuremath{-}0.340 [\ensuremath{-}0.434, \ensuremath{-}0.252] & \ding{51} \\
Centre offset, difference to fallback & CI hi $< 0$ & \ensuremath{-}0.134 [\ensuremath{-}0.191, \ensuremath{-}0.086] & \ding{51} & \ensuremath{-}0.053 [\ensuremath{-}0.087, \ensuremath{-}0.020] & \ding{51} \\
Drift, difference to fallback & CI hi $< 0$ & \ensuremath{-}0.151 [\ensuremath{-}0.369, \ensuremath{-}0.028] & \ding{51} & \ensuremath{-}0.008 [\ensuremath{-}0.018, \ensuremath{-}0.002] & \ding{51} \\
Motion error, difference to fallback & CI hi $< 0$ & \ensuremath{-}0.073 [\ensuremath{-}0.129, \ensuremath{-}0.033] & \ding{51} & \ensuremath{-}0.014 [\ensuremath{-}0.026, \ensuremath{-}0.006] & \ding{51} \\
Median text-height ratio & $[0.8, 1.2]$ & 0.730 & \textbf{\ding{55}} & 0.965 & \ding{51} \\
\multicolumn{6}{l}{\textit{Containment in the carrier}}\\
Mean containment (fallback) & $\geq 0.90$ & 0.989 (0.971) & \ding{51} & 0.945 (0.917) & \ding{51} \\
Containment, difference to fallback & CI lo $\geq -0.03$ & +0.018 [+0.007, +0.030] & \ding{51} & +0.028 [+0.011, +0.045] & \ding{51} \\
Severe-failure audit; layout acceptability & human & 3 automatic candidates & -- & 12 automatic candidates & -- \\
\bottomrule
\end{tabular}}
\end{table*}


\subsection{The sealed cross-generator test}
\label{sec:supp_g4_t4}
\textbf{Protocol.} T4 was generated after the freeze, and its clips were classified and processed only by the frozen code. Before each generator's run, a content-based isolation audit confirmed that no T4 target string occurs in the development data, that the development folders hold only development prompts, and that no method output existed for T4; the full pipeline then ran once per generator and the verdict was computed once. Of the 480 clips per generator, CogVideoX-5B gave \gTCogLone{} clips with readable matches, \gTCogLtwo{} with pseudo-text and \gTCogLthree{} without a text region, Wan2.2-TI2V-5B \gTWanLone, \gTWanLtwo{} and \gTWanLthree. The synthetic subsets hold \gTCogN{} and \gTWanN{} clips (\gTCogPrompts{} and \gTWanPrompts{} prompts), the natural subsets \gTCogNatAll{} and \gTWanNatAll{} (\gTCogNatN{} and \gTWanNatN{} with a deployed \method output). Every clip is evaluated, including the tracker's localisation failures (\gTCogLocFail\% and \gTWanLocFail\%); no synthetic clip fell back to \method.

\textbf{Results} (\cref{tab:g4}). The branch reads \gTCogRead\% and \gTWanRead\% of the synthetic clips against \gTCogReadFb\% and \gTWanReadFb\% for the static fallback and \gTCogReadVthree\% and \gTWanReadVthree\% for v3; the pooled gain is \gTReadGain{} points [\gTReadGainLo, \gTReadGainHi] (CogVideoX \gTCogReadGain, Wan2.2 \gTWanReadGain). On natural pseudo-text it reads \gTCogNatRead\% and \gTWanNatRead\% against \gTCogNatBone\% and \gTWanNatBone\% for the deployed \method (\gTCogNatGain{} [\gTCogNatGainLo, \gTCogNatGainHi] and \gTWanNatGain{} [\gTWanNatGainLo, \gTWanNatGainHi] points) and breaks \gTCogNatBroken{} and \gTWanNatBroken{} originally readable clips where \method breaks \gTCogNatBoneBroken{} and \gTWanNatBoneBroken. Its background stays within \gTBgMaxGap\,dB of the identity re-encoding in both subsets of both generators; the deployed \method is lower on the natural clips (\gTCogBgNatBone{} and \gTWanBgNatBone\,dB) because its detection-box layout is not restricted to the allowed region used here. With identical rendering, the new tracker lowers the raw jitter relative to the trajectory selection of v2/v3 (\gTCogCtrlJit{} and \gTWanCtrlJit), and it places the text closer to the truth than the fallback in offset, drift and motion error on both generators. Mean containment is \gTCogCont{} and \gTWanCont{} (fallback \gTCogContFb{} and \gTWanContFb); the original 95\% standard is reached on CogVideoX and not on Wan2.2. Three registered gates fail: on CogVideoX the median height ratio is \gTCogHeight{} (D4: \gDCogHeight) and the CI of the normalised jitter difference reaches \gTCogJitNHi; on Wan2.2 the CI of the normalised flicker difference reaches \gTWanFlkNHi. The raw jitter and flicker are higher than the fallback's on both generators (CogVideoX \gTCogJitRaw{} and \gTCogFlkRaw, Wan2.2 \gTWanJitRaw{} and \gTWanFlkRaw), as expected from the higher-contrast rendered text, which is why the normalised measures are the registered gates. The automatic rules flag \gTCogSevere{} and \gTWanSevere{} synthetic clips as severe-failure candidates, all by the crossing rule (containment below 0.80) and none by clipping or missing characters; they are candidates, not confirmed failures.

\textbf{Verdict.} \gTGatesFailed{} of the \gTGatesTotal{} automatic gates fail, so the registered verdict is a failure whatever the two human gates, still pending on both generators, would show. By the decision table fixed before the test, v4.1 is not merged into \method, its results are reported per generator as measured, and it is not modified after the unsealing. Glyph-free reconstruction remains an exploratory extension, and \method (B1) remains the method of this paper.

\subsection{Diagnostic human audit}
\label{sec:supp_g4_audit}
Because the human gates could no longer change the verdict, the registered severe-failure audit was replaced by a smaller diagnostic audit, fixed before any rating: every automatic candidate (\gTCogSevere{} CogVideoX, \gTWanSevere{} Wan2.2) and 30 clips per generator from the unflagged pool, stratified by prompt with a fixed seed, 75 clips in total. Each clip is judged by exactly 3 of 12 independent raters (225 judgements) on one page that shows the input with the allowed region, the edited clip and the target, and shows neither the method nor the candidate status. Question~1 asks whether a severe geometric problem occurs---text on the wrong carrier or crossing carriers, severe clipping or overflow, covering adjacent objects, damage to the carrier's structure; question~2 asks whether the layout is plausible and matches the carrier's geometry. Font aesthetics and material are explicitly not severe. A clip counts as severe (layout acceptable) when at least 2 of its 3 raters say so; raw labels and Fleiss' $\kappa$ are reported. The audit describes what the automatic candidates are and gives a small random-sample picture; it cannot establish a population rate of severe failures, and a candidate the raters do not confirm is not counted as a failure. \pending{Results are added when the ratings are complete.}
